\section{Discussion, Threats to Validity, and Limitations}
\label{sec:8}

\subsection{Discussion and Practical Implications}
\label{sec:8.1}

\paragraph{When to use Topo-QoS, and when to use HGT-QoS}
The results do not support a simple recommendation of the learned model over the closed-form one, and we set out the trade-off as measured rather than as hoped.

\begin{enumerate}
    \item \textbf{Training-free ranking (\texttt{Topo-QoS}).} It requires no training, no checkpoint storage, and no retraining as a corpus evolves, reaching $\rho = 0.553$ zero-shot across twelve unseen synthetic architectures and $0.526$ across five open-source systems. Nothing in this study establishes that a learned model beats it on ranking: the margin is $+0.085$ ($p = 0.151$) with an interval spanning zero, and on the one real system where the two diverge sharply it is the heuristic that wins ($0.888$ vs.\ $0.649$ on Cloud Microservices). For a team that wants a criticality ordering and nothing more, this is the defensible default, and we say so despite proposing the alternative.
    \item \textbf{One learned mechanism, not two (\texttt{GAT-N-QoS} or \texttt{HGT}).} The clearest learned result is that either relation typing or a QoS edge channel lifts an untyped, unweighted model substantially --- $+0.234$ and $+0.287$ respectively, both surviving Holm correction --- while adding the second to the first does not. A practitioner should therefore pick one. On cost grounds the untyped QoS-weighted model is the better pick: $\rho = 0.604$ at $28{,}168$ parameters against HGT-QoS's $0.638$ at $434{,}620$. Note also that the untyped, \emph{unweighted} model (GAT-N, $\rho = 0.317$) is worse than the training-free heuristic, so ``use a GNN'' is not by itself sound advice here; which relational signal it is given decides whether it is worth training at all.
    \item \textbf{Capabilities without a closed-form counterpart.} Two remain, and neither is evaluated here. Typed relational attention exposes \emph{which} channels mediate a cascade rather than only which components rank highly (Supplementary Section~S8 illustrates this on one topology and is explicitly not evidence of a general effect). Relationship-level criticality ($I_{\text{edge}}$, Eq.~\eqref{eq:edge_crit}) scores individual dependencies rather than components, which is what circuit-breaker or bulkhead placement actually requires and which a node ranking cannot express; this paper defines the edge oracle and the model's edge head but presents no evaluation of one against the other, so the capability is available and untested. Where the extra machinery is affordable these are the substantive reasons to prefer a typed model --- as design arguments, not as measured advantages.
\end{enumerate}

\paragraph{A gate we can no longer recommend}
An earlier version of this work proposed a tiered gate: run the learned model by default and fall back to \texttt{Topo-QoS} when the model's prediction dispersion $\hat{\sigma}$ fell below a threshold. That heuristic does not replicate: across twelve folds, $\hat{\sigma}$ correlates with the margin over \texttt{Topo-QoS} at $\rho_s = -0.126$ for HGT-QoS (the wrong sign, Section~\ref{sec:7.2.3}). We therefore withdraw the automated fallback recommendation. Instead, because both engines run in seconds, our results support running both concurrently and escalating ranking disagreements to human architectural review.

\paragraph{Dual-Engine Consensus Protocol}
Because both engines run in seconds, SaG ships a dual mode (\texttt{saag-predict} with \texttt{-{}-predictor-mode dual}) that scores each manifest with \texttt{HGT-QoS} and \texttt{Topo-QoS} together and partitions the result: components in the top-$K$ set of both are unanimous risks worth prioritising, while components the two rank far apart are escalated to human architectural review rather than resolved automatically. The design follows from the measurement --- we have no label-free signal that says which engine to trust on an unseen architecture, so the honest response is to surface the disagreement instead of hiding it behind a threshold. We report the protocol as a deployment recommendation; its triage value is not evaluated here.

\paragraph{Role of the Explanation Layer}
The RM attribution profile ($Q(v)$, Section~\ref{sec:5}) provides transparent, standards-compliant architectural diagnostics aligned with ISO/IEC 25010. By separating single-point-of-failure exposure (Availability) from wide error propagation reach (Fault Tolerance), RM provides qualitative remediation guidance (e.g., distinguishing whether a component requires replication or decoupling) that purely numeric rankers and simulation oracles cannot provide.

\subsection{Performance and Computational Sustainability Implications}
\label{sec:8.2}

\paragraph{What sustainability means for a pre-deployment gate}
Green software engineering distinguishes the energy a system consumes from the energy its \emph{development and assurance} consume \cite{calero2015green}, and the machine-learning literature has concentrated on the second \cite{schwartz2020greenai,strubell2019energy,patterson2021carbon,georgiou2022greenframeworks,lannelongue2021green,verdecchia2023greenai}. Reliability assurance sits in the same category and is rarely measured: chaos engineering, hardware-in-the-loop benches and staging-cluster fault injection consume cluster-hours per sweep, on every pull request that triggers them. Two findings bear on this, at the strength the measurements support.

The first is unambiguous. \emph{Within} the pipeline the neural model is negligible: a $56\,\text{ms}$ forward pass on a 2{,}000-component system against $239\,\text{s}$ for deterministic structural feature extraction, a ratio of $4{,}259\times$ (Section~\ref{sec:7.5}). Whatever pre-deployment dependability analysis costs, adding a graph neural network is not what makes it expensive --- which, for a special issue asking whether AI techniques can be afforded in a sustainability-conscious process, is the directly relevant result.

The second is a difference in kind rather than degree, and it is weaker than it first appears. What manifest-time analysis eliminates is \emph{infrastructure}: nothing must be deployed, kept warm, or torn down, and no fault is injected into anything a user could be holding. But that argument does not distinguish SaG from its own oracle, which also reads a manifest and costs eleven times less (Section~\ref{sec:7.5.oracle}); it distinguishes static analysis in general from chaos engineering on a provisioned cluster, and we have not measured a chaos-engineering baseline, so its magnitude is an assertion rather than a result.

\paragraph{The efficiency claim we withdraw, and where the cost actually sits}
Static analysis does not reduce computation relative to in-process simulation: global connectivity degradation ($82.7\,\text{s}$ on the Enterprise mesh) is roughly eleven times slower than BFS cascade traversal ($7.2\,\text{s}$). We withdraw any general claim of computational efficiency relative to simulation. The expense is concentrated in one metric and is a deliberate accuracy purchase: the Connectivity Degradation Index is computed for every node in the main component rather than for articulation points alone, and gating it would restore roughly an order of magnitude of speed at the price of a degenerate Availability score. A deployment needing the speed more than the single-point-of-failure sensitivity could make the opposite trade. This also locates where optimization effort belongs: the sustainability of this framework is a question about one graph-theoretic routine, not about its use of machine learning. Settling it comprehensively requires energy counters rather than wall-clock time \cite{lannelongue2021green}; and while preventing cascading failures plausibly reduces datacenter compute lost to retry storms and restart loops, we treat that as motivation rather than an empirical claim.

\subsection{Threats to Validity}
\label{sec:8.3}

\paragraph{Construct Validity} Our primary ground-truth impact oracle $I^*(v)$ is derived from discrete-event cascade simulation on structural models rather than live outages. Evaluating construct divergence against the queue-flow discrete-event simulator $I_{\text{dyn}}(v)$ and composite oracle $I_{\text{comp}}(v)$ (Supplementary Section~S9) reveals substantial but sub-ceiling correlation with $I_{\text{dyn}}$ ($\rho = 0.620$ against a $0.811$--$1.000$ label test--retest ceiling), which traverses related but not identical failure semantics. However, agreement with $I_{\text{comp}}$ is moderate ($\rho = 0.395$), and top-$K$ critical-set Jaccard reaches only $0.27$--$0.37$ due to non-linear cascade threshold sensitivity and tied zero-inflation (Supplementary Section~S2). Thus, convergent validity across simulation paradigms holds for ranking but is weaker for critical-set boundaries; no oracle is validated against live outages (Section~\ref{sec:4.4}). For Maintainability, $I_M(v)$ traverses the same dependency topology from which $M(v)$ is scored; independent churn validation remains future work.

\paragraph{Internal Validity} Potential feature leakage is prevented by strict graph view separation: predictors operate exclusively on $G_{\text{analysis}}$, whereas ground-truth simulation oracles operate on $G_{\text{structural}}$, formally asserted in continuous integration. Substrate parity is rigorously maintained: learned models (HGT-QoS, GAT-N-QoS) share identical training sets, depths, and early-stopping rule---a validation split within the primary training graph (\S\ref{sec:6.3}), applied uniformly across variants. In the QoS schema, four dimensions (reliability, durability, transport priority, heterogeneity flag) capture active operational middleware configurations, and two further dimensions encode the declared deadline, which is populated for $75\%$ of topics across the twelve evaluation scenarios. The remaining dimension (max blocking) is a reserved extension point that is zero throughout the corpus. 

\paragraph{External Validity} Our evaluation spans twelve synthetic architectures and five open-source systems. Zero-shot transfer to real systems is \emph{not} demonstrated on active components (mean $\rho_{>0} = +0.265$, inverting on the two microservice call trees, Section~\ref{sec:7.4.learned}). The explanation we offer for that pattern --- that pub-sub meshes and synchronous RPC call trees propagate failure in opposite directions along the edges the model learned --- was arrived at after seeing which systems failed, on a sample of five. It is a conjecture consistent with the data, not a tested hypothesis, and three systems against two cannot support a statistical claim either way. We state it because it is actionable (it predicts which architectures the current corpus cannot prepare a model for) and label it as post-hoc for the same reason. Furthermore, while synthetic QoS profiles exhibit genuine variance (modal shares $29$--$89\%$, Section~\ref{sec:7.3.1}), their alignment with production distributions remains unverified. Finally, our timing evaluations scale to 2{,}000 components; architectures an order of magnitude larger cannot be assumed to fit PR budgets given the $O(|V|^2 + |V||E|)$ dominant stage.

\paragraph{Conclusion Validity} Given heavy-tailed impact distributions, statistical analyses use non-parametric rank correlation (Spearman $\rho$, Kendall $\tau$), bootstrap confidence intervals ($B = 2{,}000$), and paired Wilcoxon signed-rank tests, with the fold or scenario as the unit of analysis. Two hazards recur and shape how we report. Pooling across heterogeneous entity types triggers Simpson's paradox --- pooled $\rho = 0.098$ sits below every per-type value it aggregates ($0.119$--$0.566$) --- so all headline figures are stratified on a single population. And rank correlation over zero-inflated labels conflates ordering the active components with separating them from inert ones, which is why we report zero-excluded correlations alongside full-population ones wherever the label distribution permits. Where the two disagree, as on the real-world systems, we read the zero-excluded figure as the ranking result.

Two further limits bound every $p$-value and interval we report, and neither is removable by analysis. The folds are not independent replicates: any two LOSO models share ten of their eleven training graphs, so a signed-rank test over folds treats as independent observations that are strongly coupled, and there is no unbiased estimator of the variance of a cross-validation estimate under this design. Our intervals and $p$-values are therefore optimistic by an amount we cannot quantify, and we report them as the conventional summary rather than as calibrated ones. Separately, the twelve synthetic topologies are draws from a single parameterized generator sharing fan-out, colocation and QoS-categorical distributions (Supplementary Section~S5), so ``generalizes to unseen architectures'' means, strictly, unseen draws from one generator. The five open-source systems are the only architectures in this study whose structure we did not specify, and they are the ones on which transfer is not established.

\subsection{Limitations and Future Work}
\label{sec:8.4}

\paragraph{A baseline that was missing for the wrong reason} Earlier versions of this work omitted \texttt{Topo-QoS} from Table~\ref{tab:9b}, reporting that the open-source adapters carried no QoS contracts to weight with. That was incorrect. Every topic in all five adapters declares durability, reliability and transport priority, and projecting $w(t)$ onto the structural edges yields non-unit weights on roughly half of them; the obstacle was a defect in the projection guard, which computed QoS-weighted betweenness on the raw multigraph --- where Application nodes never route messages and the score is degenerate by construction --- rather than on the \texttt{DEPENDS\_ON} projection every other topological baseline uses. With that corrected the column is scorable, and Table~\ref{tab:9b} reports it. We record the error rather than quietly filling the gap, because for two revisions it meant the real-world comparison was made against the weaker of the two training-free references while the stronger one was described as unavailable.


\paragraph{The explanation layer is not validated as an explanation} SaG separates Availability from Fault Tolerance on the hypothesis of distinct repairs, but we do not evaluate whether practitioners find this actionable. Furthermore, elicited AHP weights perform worse than a uniform prior at ranking (Section~\ref{sec:7.3}); counterfactual mutation tests and user studies are needed for validation.

\paragraph{Two confounds in the typing result remain uncontrolled} The contrasts of Table~\ref{tab:contrasts} hold substrate, training set, depth and selection rule constant, but two factors separating the typed and untyped architectures are unmatched as published. \emph{Parameter budget:} on the relation set of the LOSO primary training graph, HGT-QoS carries $434{,}620$ parameters against GAT-N-QoS's $28{,}168$. \emph{Directionality:} HGT applies a reverse-direction \texttt{HGTConv} ($103{,}725$ parameters) so a node sees its upstream neighbourhood, whereas the homogeneous baseline propagates along native edge direction only --- which matters because $I^*(v)$ is a downstream-reachability functional, so a model able to look upstream has an advantage on this target unrelated to typing. The arms that would isolate each (a capacity-matched homogeneous baseline, one reading the full sixteen-dimensional edge channel, and a forward-only heterogeneous model) are implemented and registered in the replication package but were not run at the reported budget, so no value is reported for them. The consequence is specific and we do not minimise it: the surviving $+0.234$ is consistent with relational typing acting as an inductive bias, and equally consistent with a capacity or directionality advantage. Running those three arms is the single most informative extension to this study.

\paragraph{Model selection is made on the training distribution} Early stopping in LOSO uses an inner validation split within the primary training graph (Section~\ref{sec:6.3}). While a protocol limitation under distribution shift, the rule is uniform across variants and cannot manufacture the typed-versus-untyped margin. Held-out scenario validation (\texttt{-{}-inner-val-scenario auto}) is our prioritized extension.

\paragraph{No label-free reliability signal} Prediction dispersion does not replicate as a fallback indicator on this corpus (Section~\ref{sec:7.2.3}); discovering a reliable OOD confidence signal remains an open problem.

\paragraph{Scale is bounded by the deterministic stage, not the model} Our timings reach 2{,}000 components. Beyond that the binding constraint is the $O(|V|^2 + |V||E|)$ feature-extraction stage rather than the network (Section~\ref{sec:7.5}), so the useful optimizations are graph-theoretic: caching structural metrics across pull requests and recomputing only over the $k$-hop neighbourhood a manifest change touches. Training, separately, would need mini-batch subgraph sampling (GraphSAINT \cite{zeng2020graphsaint} or HGT's own layer-dependent importance sampling) to reach the $|V| > 10^5$ regime of cross-organizational service fleets, since full-batch message passing expands its receptive field exponentially with depth. Neither is a limitation this evaluation encountered; both are what a deployment at an order of magnitude more scale would encounter first.

\paragraph{Future Directions: Distributed AI, Power Testbeds, and Self-Healing}
We envision three primary extensions: (1)~modeling distributed LLM serving clusters (e.g., vLLM, DeepSpeed) to mitigate straggler-induced GPU dissipation; (2)~measuring hardware energy directly via RAPL/NVML to compare static gating against live chaos sweeps in joules; and (3)~advancing from predictive diagnostics to prescriptive synthesis, automatically generating pull requests with circuit breakers, broker replicas, and tuned QoS parameters to resolve single points of failure.
